\labeledsubsubsection{Defining a plugin}
The previous sections already prepared everything to load a plugin.
However, it is not yet defined what a \textit{plugin} actually is.

There are multiple ways of defining a \textit{plugin}.
One way would be to use specific class and method names.
Another way is to utilize \underline{interfaces} and/or \underline{abstract classes}.
In \texttt{Java}, there is a third good option: 
Using \underline{annotations} (\underline{@interfaces}).

The experiment uses an \underline{annotation} to denote a class as a \textit{plugin-class}.
A \textit{plugin-class} is the entry-point of a plugin.
It defines the name and version of a plugin and the class that is annotated will be searched for \textit{plugin-methods}, which will be handled next.
The following code shows how the class \underline{annotation} is defined in the experiment:
\lstinputlisting[language=Java, linerange={8-14}]{./code/plugin-api/src/main/java/nl/fontys/sear/plugin_system/api/PluginInfo.java}

Another \underline{annotation} is used to denote the exact \textit{plugin-method}.
Said method is the main reason why the \textit{plugin-system} wanted to load the \textit{plugin} in the first place.
In a real scenario this method would either initialize the plugin or perform a certain task that is wanted by the \textit{plugin-system}.
This is also the main attack point in the later analysis.
\lstinputlisting[language=Java, linerange={8-11}]{./code/plugin-api/src/main/java/nl/fontys/sear/plugin_system/api/Task.java}

\clearpage

An example plugin entry-point class would look like this:
\lstinputlisting[language=Java, linerange={8-9, 20-23, 24}]{./code/plugin/src/main/java/nl/fontys/sear/plugin_system/plugin/PluginMain.java}